\documentclass[aspectratio=169,11pt,t]{beamer}
\input{../../../theme/ksa-cs-v2.tex}
\usepackage{fontawesome5}
\definecolor{ksagreen}{HTML}{2A9D5C}
\newcommand{\paperlink}[1]{\href{#1}{\textcolor{ksablue}{\faLink}}}
% GDA 차시 (2026-11-27): 앤드류 응 CS229 강의 노트 Part IV (Generative Learning Algorithms) 의 GDA 를 덱 스타일로.
% 숫자는 data.py (assert), 그림은 make_figs.py. 데이터: 붓꽃 versicolor · virginica 100송이.
\setbeamerfont{author}{size=\normalsize}
\setbeamerfont{institute}{size=\small}
\graphicspath{{figs/}{../../../assets/}}

\renewcommand{\ksacredit}{Machine Learning and Deep Learning \,\textperiodcentered\, GDA}
\title{가우시안 판별 분석 (GDA)}
\subtitle{경계 대신 ``종마다 어떻게 생겼나''를 배운다: 생성 모델}
\author{Machine Learning and Deep Learning}
\institute{KSA 2027-Spring}
\date{Sumin Han (suminhan@ksa.hs.kr)}
\begin{document}
{\setbeamertemplate{footline}{}\begin{frame}[c,noframenumbering]\titlepage\end{frame}}

\section{질문}
% --- 오늘의 질문
\begin{frame}{오늘의 질문: 꽃잎 4.9 cm 는 versicolor? virginica?}
\vfill
\centering
\renewcommand{\arraystretch}{1.7}
{\Large\begin{tabular}{c|c}
\textbf{판별 모델} & \textbf{생성 모델}\\ \hline
경계를 \kb{바로} 배운다 & 종마다 \kb{어떻게 생겼나}를 배운다\\
$P(y \mid x)$ & $P(x \mid y)$ \;와\; $P(y)$ \;$\to$\; 베이즈\\
로지스틱회귀 · 소프트맥스 & \kb{GDA} · 나이브 베이즈
\end{tabular}}\\[1.0em]
{\large 코끼리와 개를 가를 때: 둘 사이에 선을 긋거나, 코끼리 · 개를 따로 그려 보고 더 닮은 쪽}\\[0.2em]
{\small GDA = Gaussian Discriminant Analysis (가우시안 판별 분석) \quad 붓꽃 두 종 50송이씩}
\vfill
\note{꽃잎 길이가 4.9센티인 붓꽃은 versicolor일까요, virginica일까요? 지금까지 본 로지스틱회귀나 소프트맥스는 두 종 사이의 경계를 바로 배웠습니다. x가 주어졌을 때 y의 확률을 곧장 모델로 만든 것이고, 이런 모델을 판별 모델이라고 부릅니다. 오늘은 반대 방향입니다. 종마다 꽃이 어떻게 생겼는지, 즉 y가 주어졌을 때 x의 분포와, 종 자체의 비율을 배우고, 베이즈 정리로 뒤집어서 분류합니다. 이런 모델을 생성 모델이라고 합니다. 앤드류 응의 예로 말하면, 코끼리와 개를 구분할 때 둘 사이에 선을 긋는 것이 판별 모델, 코끼리는 이렇게 생겼고 개는 이렇게 생겼다를 따로 배워 두고 새 동물이 어느 쪽을 더 닮았는지 보는 것이 생성 모델입니다. 지난 시간의 나이브 베이즈도 생성 모델이었습니다. 오늘은 각 종의 모양을 가우시안으로 배우는 GDA, 가우시안 판별 분석입니다.}
\end{frame}

\section{1차원}
% --- 종마다 가우시안
\begin{frame}{종마다 꽃잎 길이를 가우시안으로}
\vfill
\centering\includegraphics[width=0.58\linewidth]{gda_1d.png}\par\vskip0.2em
{\large versicolor: 평균 \textcolor{ksagreen}{4.26} cm \quad virginica: 평균 \textcolor{ksablue}{5.55} cm \quad 표준편차는 둘이 같게 \kb{0.51} cm}\\[0.2em]
{\large 4.9 cm 에서 종의 높이 (가능도): \;\textcolor{ksagreen}{0.355} vs \textcolor{ksablue}{0.344} \;$\to$\; 거의 같다}\\[0.2em]
{\small 표준편차를 종마다 따로 재면 0.47 · 0.55 \quad 같게 두는 이유는 뒤에서}
\vfill
\note{먼저 종마다 꽃잎 길이가 어떻게 생겼는지 봅니다. 연한 막대가 실제 50송이씩의 히스토그램이고, 굵은 곡선이 거기에 맞춘 가우시안입니다. versicolor는 평균 4.26센티, virginica는 평균 5.55센티입니다. 표준편차는 두 종이 같다고 두고, 두 종을 합쳐 잰 0.51센티를 씁니다. 따로 재면 0.47과 0.55로 비슷합니다. 왜 같게 두는지는 조금 뒤에 나옵니다. 이제 4.9센티에서 두 종의 높이를 읽습니다. versicolor는 0.355, virginica는 0.344입니다. 이 높이가 가능도, 그 종이라면 4.9센티 꽃이 나올 정도입니다. 4.9센티는 두 곡선이 거의 만나는 곳이라 높이도 거의 같습니다.}
\end{frame}

% --- 베이즈로 확률
\begin{frame}{베이즈 정리로 뒤집는다}
\vfill
\centering
{\Large\[ \underbrace{P(\text{virginica}\mid 4.9)}_{\text{posterior}} = \frac{\textcolor{ksablue}{0.344}\times\textcolor{ksared}{0.5}}{\underbrace{\textcolor{ksablue}{0.344}\times\textcolor{ksared}{0.5} + \textcolor{ksagreen}{0.355}\times\textcolor{ksared}{0.5}}_{\text{evidence}}} = \kb{0.49} \]}\vskip0.4em
\renewcommand{\arraystretch}{1.25}
{\normalsize\begin{tabular}{lll}
\textcolor{ksablue}{0.344}, \textcolor{ksagreen}{0.355} & 가능도 likelihood $P(x \mid y)$ & 각 종이라면 4.9 cm 가 나올 정도 (종의 높이)\\
\textcolor{ksared}{0.5} & 사전확률 prior $P(y)$ & 꽃을 보기 전 각 종일 확률 (50송이씩)\\
0.49 & 사후확률 posterior $P(y \mid x)$ & 4.9 cm 를 본 뒤 virginica 일 확률
\end{tabular}}
\vfill
\note{나이브 베이즈에서 본 베이즈 정리를 그대로 씁니다. 분자는 virginica의 가능도 0.344 곱하기 사전확률 0.5, 분모는 두 종의 경우를 더한 것입니다. 계산하면 0.49, 4.9센티 꽃이 virginica일 확률은 49퍼센트입니다. 사전확률은 꽃을 보기 전에 각 종일 확률인데, 두 종이 50송이씩이라 0.5입니다. 나이브 베이즈와 다른 점은 가능도를 어떻게 구하느냐 하나입니다. 나이브 베이즈는 단어를 세서 구했고, GDA는 가우시안 종의 높이로 구합니다.}
\end{frame}

% --- 사후확률 곡선
\begin{frame}{길이마다 사후확률을 그리면: S 자}
\vfill
\centering\includegraphics[width=0.58\linewidth]{gda_post.png}\par\vskip0.2em
{\large GDA 로 구한 확률도 \kb{로지스틱 곡선} 모양 \quad 그런데 4.9 cm 에서 \textcolor{ksablue}{49\%} vs \textcolor{ksagolddk}{58\%}}\\[0.2em]
{\large 같은 꼴, \kb{다른 숫자}: 왜 S 자인지, 왜 다른지}
\vfill
\note{4.9센티만이 아니라 모든 길이에서 사후확률을 계산해 그려 봅니다. 파란 곡선이 GDA입니다. 짧으면 0, 길면 1로 가는 S 자, 로지스틱 곡선 모양입니다. 같은 100송이로 로지스틱회귀를 경사하강법으로 학습한 것이 금색 점선입니다. 둘 다 S 자인데 기울기가 다르고, 4.9센티에서 GDA는 49퍼센트, 로지스틱회귀는 58퍼센트입니다. 두 가지를 알아봅시다. 왜 GDA도 S 자가 되는지, 그리고 같은 꼴인데 왜 숫자가 다른지.}
\end{frame}

% --- 유도
\begin{frame}{왜 S 자인가: 로그 오즈가 직선이 된다}
\vfill
\centering
{\large\begin{align*}
\log\frac{P(y{=}1\mid x)}{P(y{=}0\mid x)} &= \log\frac{\phi\,\mathcal{N}(x;\mu_1,\sigma^2)}{(1-\phi)\,\mathcal{N}(x;\mu_0,\sigma^2)} && \text{베이즈 (분모는 약분)}\\[0.2em]
&= -\frac{(x-\mu_1)^2}{2\sigma^2} + \frac{(x-\mu_0)^2}{2\sigma^2} + \log\frac{\phi}{1-\phi} && \text{가우시안 식, } \tfrac{1}{\sqrt{2\pi}\sigma} \text{ 약분}\\[0.2em]
&= \underbrace{\frac{\mu_1-\mu_0}{\sigma^2}}_{w_1}\,x \;+\; \underbrace{\frac{\mu_0^2-\mu_1^2}{2\sigma^2} + \log\frac{\phi}{1-\phi}}_{w_0} && \textcolor{ksared}{x^2 \text{ 항이 지워진다}}
\end{align*}}\vskip0.3em
{\Large $\to$ \;$P(y{=}1\mid x) = \mathrm{sigmoid}(w_0 + w_1 x)$ \;=\; \kb{로지스틱 꼴}}\\[0.3em]
{\small $\phi$ = virginica 비율 (0.5) \quad 표준편차를 종마다 다르게 두면 $x^2$ 가 남는다 $\to$ S 자가 아니다}
\vfill
\note{왜 S 자인지 식으로 봅시다. 로지스틱회귀 시간에 본 로그 오즈, virginica일 확률 나누기 versicolor일 확률에 로그를 씌웁니다. 베이즈 정리에서 분모는 같으니 약분되고, 파이 곱하기 virginica의 가우시안 나누기 1 빼기 파이 곱하기 versicolor의 가우시안이 남습니다. 파이는 virginica의 비율입니다. 가우시안 식을 넣으면 앞의 1 나누기 루트 2파이 시그마는 두 종이 같아서 약분되고, 지수 부분만 남습니다. 제곱을 전개하면 x제곱 항이 두 개 나오는데, 시그마가 같으니 서로 지워집니다. 그러면 x의 일차식, w1 x 더하기 w0만 남습니다. 로그 오즈가 x의 직선이면 확률은 시그모이드 w0 더하기 w1 x, 로지스틱회귀와 똑같은 꼴입니다. 표준편차를 같게 둔 이유가 여기 있습니다. 종마다 다르게 두면 x제곱 항이 남아서 S 자가 아닌 다른 모양이 됩니다.}
\end{frame}

% --- 같은 꼴 다른 w
\begin{frame}{같은 꼴, 다른 숫자: $w$ 를 구하는 방법이 다르다}
\vfill
\centering
\renewcommand{\arraystretch}{1.7}
{\Large\begin{tabular}{l|c|c|c}
 & $\mathrm{sigmoid}(w_0 + w_1 x)$ & 경계 & 4.9 cm\\ \hline
GDA & $w_0 = -24.6, \; w_1 = 5.02$ & 4.91 cm & 49\%\\
로지스틱회귀 & $w_0 = -43.8, \; w_1 = 9.00$ & 4.86 cm & 58\%
\end{tabular}}\\[1.0em]
{\large GDA: \kb{평균 · 표준편차를 재서} $w_1 = \dfrac{5.55 - 4.26}{0.51^2} = 5.02$ \quad (경사하강법 없음)}\\[0.4em]
{\large 로지스틱회귀: 경계 근처 점을 맞히도록 \kb{경사하강법}}
\vfill
\note{두 모델 모두 시그모이드 w0 더하기 w1 x이지만 숫자가 다릅니다. GDA는 w0 마이너스 24.6, w1 5.02, 경계 4.91센티. 로지스틱회귀는 w0 마이너스 43.8, w1 9.00, 경계 4.86센티입니다. 차이는 w를 구하는 방법입니다. GDA는 앞의 식대로 평균과 표준편차를 재서 바로 계산합니다. w1은 5.55 빼기 4.26을 0.51 제곱으로 나눈 5.02입니다. 경사하강법이 필요 없습니다. 로지스틱회귀는 경사하강법으로, 실제로 맞히고 틀리는 경계 근처의 꽃에 맞춰 w를 고칩니다. 그래서 로지스틱회귀는 경계 근처가 더 가파릅니다. 같은 꼴의 모델을 두 가지 다른 생각으로 구한 것입니다.}
\end{frame}

\section{2차원}
% --- 다변량 가우시안
\begin{frame}{특징이 둘이면: 다변량 가우시안}
\vfill
\centering
{\LARGE\[ \mathcal{N}(\mathbf{x};\,\boldsymbol{\mu},\,\Sigma) = \frac{1}{2\pi\,|\Sigma|^{1/2}}\exp\!\Big(-\tfrac12\,(\mathbf{x}-\boldsymbol{\mu})^{\!\top}\,\Sigma^{-1}\,(\mathbf{x}-\boldsymbol{\mu})\Big) \]}\vskip0.5em
\renewcommand{\arraystretch}{1.4}
{\large\begin{tabular}{lll}
$\mathbf{x}$ & (꽃잎 길이, 꽃잎 너비) & 점 하나\\
$\boldsymbol{\mu}$ & 평균 벡터 & 종의 가운데\\
$\Sigma$ & 공분산 행렬 ($2\times 2$) & 종의 \kb{폭과 기울기}\\
$|\Sigma|$, $\Sigma^{-1}$ & 행렬식, 역행렬 & 1차원의 $\sigma^2$, $1/\sigma^2$ 자리
\end{tabular}}
\vfill
\note{특징을 꽃잎 길이와 꽃잎 너비 두 개로 늘려 봅시다. 그러면 가우시안도 2차원이 됩니다. 식은 복잡해 보이지만 1차원과 같은 모양입니다. 1차원의 x 빼기 뮤의 제곱 나누기 시그마 제곱이, 벡터 x 빼기 뮤에 공분산 행렬의 역행렬을 앞뒤로 곱한 것으로 바뀌었고, 앞의 시그마가 공분산 행렬의 행렬식으로 바뀌었습니다. 평균 벡터 뮤는 종의 가운데, 공분산 행렬 시그마는 종이 얼마나 넓게, 어느 방향으로 퍼졌는지를 담습니다. 공분산 행렬을 읽는 법을 다음 쪽에서 봅니다.}
\end{frame}

% --- Σ 읽기
\begin{frame}{공분산 행렬 $\Sigma$ 읽기}
\vfill
\centering\includegraphics[width=0.86\linewidth]{gda_sigma.png}\par\vskip0.3em
{\large 대각 $=$ 각 특징의 \kb{분산} (얼마나 넓게) \quad 비대각 $=$ 두 특징이 \kb{함께 움직이는 정도}}\\[0.2em]
{\small 타원 = 높이가 같은 곳 (등고선), 안쪽 $1\sigma$ · 바깥 $2\sigma$}
\vfill
\note{공분산 행렬은 2 곱하기 2 행렬입니다. 왼쪽처럼 대각이 1, 1이고 나머지가 0이면 둥근 종입니다. 가운데처럼 대각의 첫 번째가 2.5로 크면 가로로 넓게 퍼집니다. 대각은 각 특징의 분산, 얼마나 넓게 퍼졌는지입니다. 오른쪽처럼 비대각이 0.8이면 비스듬한 타원이 됩니다. 한 특징이 크면 다른 특징도 큰 경향, 함께 움직이는 정도가 비대각입니다. 꽃잎은 길수록 넓은 경향이 있으니 붓꽃은 오른쪽처럼 비스듬할 겁니다.}
\end{frame}

% --- 2D GDA
\begin{frame}{붓꽃 100송이로 2차원 GDA}
\vfill
\begin{columns}[c]
\begin{column}{0.56\textwidth}\centering\includegraphics[width=\linewidth]{gda_2d.png}\end{column}
\begin{column}{0.42\textwidth}
{\large 평균 (\textcolor{ksagreen}{$\times$}, \textcolor{ksablue}{$\times$})}\\[0.1em]
{\normalsize versicolor $(4.26,\ 1.33)$, virginica $(5.55,\ 2.03)$}\\[0.6em]
{\large 공유 공분산} \;{\normalsize $\Sigma = \begin{bmatrix}0.257 & 0.060\\ 0.060 & 0.056\end{bmatrix}$}\\[0.6em]
{\large 같은 $\Sigma$ $\to$ 같은 모양 타원 두 개}\\[0.2em]
{\large $\to$ 경계는 \kb{직선} \;(정확도 94\%)}\\[0.4em]
{\small $P(y{=}1 \mid \mathbf{x}) = \mathrm{sigmoid}(w_0 + \mathbf{w}\cdot\mathbf{x})$, \;$\mathbf{w} = \Sigma^{-1}(\boldsymbol{\mu}_1 - \boldsymbol{\mu}_0)$}
\end{column}
\end{columns}
\vfill
\note{이제 실제 붓꽃 100송이로 2차원 GDA를 합니다. 가위표가 두 종의 평균입니다. versicolor는 길이 4.26, 너비 1.33, virginica는 5.55, 2.03입니다. 공분산 행렬은 두 종이 하나를 같이 씁니다. 대각 0.257과 0.056은 길이와 너비의 분산이고, 비대각 0.060이 양수라 길이가 길면 너비도 넓은 비스듬한 타원입니다. 같은 공분산이니 두 종의 타원은 모양이 같고 위치만 다릅니다. 1차원과 똑같이 x제곱 항이 지워져서 경계는 직선이고, 사후확률은 시그모이드 w0 더하기 w 점곱 x, 로지스틱 꼴입니다. w는 공분산의 역행렬 곱하기 두 평균의 차이입니다. 이 경계로 100송이 중 94퍼센트를 맞힙니다.}
\end{frame}

% --- 학습 = 세기
\begin{frame}{GDA 학습 $=$ 세고, 평균 내기}
\vfill
\centering
{\Large\begin{align*}
\phi &= \frac{1}{m}\sum_{i=1}^{m} 1\{y^{(i)} = 1\} && \text{virginica 비율} = 0.5\\[0.2em]
\boldsymbol{\mu}_k &= \text{종 } k \text{ 인 꽃들의 평균} && (4.26, 1.33),\ (5.55, 2.03)\\[0.2em]
\Sigma &= \frac{1}{m}\sum_{i=1}^{m}\big(\mathbf{x}^{(i)} - \boldsymbol{\mu}_{y^{(i)}}\big)\big(\mathbf{x}^{(i)} - \boldsymbol{\mu}_{y^{(i)}}\big)^{\!\top} && \text{자기 종 평균에서 퍼진 정도}
\end{align*}}\vskip0.4em
{\large 데이터가 나올 확률을 가장 크게 (최대 가능도) 하면 이 식이 \kb{바로} 나온다 \quad 반복 없음}\\[0.2em]
{\small $1\{\cdot\}$ = 맞으면 1, 아니면 0 \quad $m$ = 꽃 수 (100)}
\vfill
\note{GDA의 학습은 아주 단순합니다. 사전확률 파이는 virginica의 비율, 0.5입니다. 평균은 종마다 꽃들의 평균입니다. 공분산은 각 꽃이 자기 종의 평균에서 얼마나 퍼졌는지를 모두 모아 평균 낸 것입니다. 학습 데이터가 나올 확률을 가장 크게 하는 최대 가능도를 풀면 이 세 식이 바로 나옵니다. 로지스틱회귀처럼 경사하강법을 여러 번 돌릴 필요가 없습니다. 세고, 평균 내면 끝입니다.}
\end{frame}

% --- QDA
\begin{frame}{$\Sigma$ 를 종마다 따로 두면: 곡선 경계}
\vfill
\centering\includegraphics[width=0.86\linewidth]{gda_qda.png}\par\vskip0.3em
{\large 같은 $\Sigma$: $x^2$ 항이 지워져 \kb{직선} (LDA) \quad 종마다 $\Sigma_k$: $x^2$ 항이 남아 \kb{곡선} (QDA)}\\[0.2em]
{\small LDA = Linear Discriminant Analysis, QDA = Quadratic $\cdots$ \quad 앤드류 응의 GDA 는 같은 $\Sigma$ (LDA)}
\vfill
\note{공분산을 두 종이 같이 쓰지 않고 종마다 따로 두면 어떨까요? 오른쪽 그림입니다. 이제 타원 모양이 종마다 다르고, x제곱 항이 지워지지 않아서 경계가 곡선이 됩니다. 이 경우를 QDA, 이차 판별 분석이라고 부르고, 같은 공분산을 쓰는 왼쪽을 LDA, 선형 판별 분석이라고 부릅니다. 앤드류 응의 강의에서 GDA라고 부르는 것은 왼쪽, 같은 공분산입니다. 곡선 경계는 이 데이터에서 97퍼센트로 조금 더 맞히지만, 숫자를 더 많이 배워야 해서 데이터가 적으면 불안정해집니다.}
\end{frame}

\section{생성 vs 판별}
% --- 비교
\begin{frame}{GDA vs 로지스틱회귀: 어느 쪽이 나은가}
\vfill
\centering
\renewcommand{\arraystretch}{1.6}
{\large\begin{tabular}{l|c|c}
 & GDA (생성) & 로지스틱회귀 (판별)\\ \hline
가정 & 종마다 \kb{가우시안} (강하다) & 확률이 S 자 (약하다)\\
학습 & 평균 · 공분산 재기 & 경사하강법\\
가정이 맞으면 & \kb{적은 데이터로도 좋다} & 데이터가 더 필요\\
가정이 틀리면 & 틀린 모양에 맞춘다 & \kb{더 튼튼하다}
\end{tabular}}\\[1.0em]
{\large 가우시안이면 $\to$ S 자 (항상) \quad S 자라고 $\to$ 가우시안은 \kq{아니다}}
\vfill
\note{그러면 어느 쪽을 쓸까요? GDA는 종마다 가우시안이라는 강한 가정을 합니다. 로지스틱회귀는 확률이 S 자라는 약한 가정만 합니다. 방금 본 것처럼 가우시안이면 확률은 언제나 S 자가 되지만, 거꾸로 S 자라고 해서 가우시안이라는 보장은 없습니다. 예를 들어 포아송 분포를 따르는 개수 데이터도 S 자를 만듭니다. 그래서 가정이 맞으면 GDA가 더 많은 것을 알고 출발하니 적은 데이터로도 좋고, 가정이 틀리면 GDA는 틀린 모양에 억지로 맞추니 로지스틱회귀가 더 튼튼합니다. 실제로 데이터로 확인해 봅시다.}
\end{frame}

% --- 적은 데이터 실험
\begin{frame}{학습 꽃이 적으면: 붓꽃으로 확인}
\vfill
\centering\includegraphics[width=0.56\linewidth]{gda_small.png}\par\vskip0.2em
{\large 꽃 10 · 20송이로 배우면 GDA 가 \kb{1.5\%p 남짓} 높다 \quad 70송이면 거의 같다}\\[0.2em]
{\large 단, 4송이면 둘 다 불안정하고 로지스틱회귀가 높다 (83.7\% vs 86.7\%)}\\[0.1em]
{\small 꽃잎 길이 · 너비 2차원, 종마다 절반씩 무작위로 뽑아 학습 · 나머지로 시험을 500번 반복한 평균}
\vfill
\note{붓꽃 100송이에서 학습에 쓸 꽃 수를 바꿔 가며 실험했습니다. 종마다 절반씩 무작위로 뽑아 학습하고 나머지로 시험하기를 500번 반복한 평균입니다. 꽃 10송이로 배우면 GDA는 92.1퍼센트, 로지스틱회귀는 90.5퍼센트, 20송이면 93.5 대 91.9로 GDA가 1.5퍼센트포인트 남짓 높습니다. 70송이쯤 되면 93.9 대 93.7로 거의 같아집니다. 붓꽃의 꽃잎 길이와 너비는 가우시안에 꽤 가까워서, 가정이 맞는 경우의 예상대로입니다. 다만 4송이, 종마다 2송이만 쓰면 둘 다 불안정하고 오히려 로지스틱회귀가 높습니다. 2송이로는 공분산을 제대로 잴 수가 없기 때문입니다. 깔끔한 법칙이 아니라 경향이라는 점을 기억해 두세요.}
\end{frame}

% --- 생성
\begin{frame}{생성 모델은 꽃을 만들 수 있다}
\vfill
\centering\includegraphics[width=0.84\linewidth]{gda_sample.png}\par\vskip0.3em
{\large 배운 가우시안에서 점을 뽑으면 \kb{가짜 꽃} \quad 판별 모델은 경계만 알아서 못 한다}\\[0.2em]
{\large 이 생각을 크게 키운 것: VAE · GAN · 확산 모델 (14주차)}
\vfill
\note{생성 모델이라는 이름에는 이유가 있습니다. GDA는 종마다 꽃이 어떻게 생겼는지 가우시안으로 배웠으니, 그 가우시안에서 점을 무작위로 뽑으면 새 꽃을 만들 수 있습니다. 왼쪽이 진짜 꽃 100송이, 오른쪽이 GDA가 만든 가짜 꽃 100송이입니다. 비스듬한 모양까지 비슷하죠. 로지스틱회귀 같은 판별 모델은 경계만 알기 때문에 이런 일을 할 수 없습니다. 이 생각, 데이터가 어떻게 생겼는지를 배워서 새것을 만든다는 생각을 아주 크게 키운 것이 14주차에 볼 VAE, GAN, 확산 모델입니다.}
\end{frame}

% --- 나이브 베이즈와 관계
\begin{frame}{나이브 베이즈와 같은 가족}
\vfill
\centering
\renewcommand{\arraystretch}{1.7}
{\large\begin{tabular}{l|c|c}
 & 나이브 베이즈 & GDA\\ \hline
특징 $x$ & 단어가 있다 / 없다 (0, 1) & 길이 · 너비 (실수)\\
$P(x \mid y)$ & 단어마다 베르누이 & 가우시안\\
특징끼리 & \kb{독립} 가정 & $\Sigma$ 로 \kb{함께 움직임}까지\\
분류 & \multicolumn{2}{c}{둘 다 $P(x \mid y)\,P(y)$ $\to$ 베이즈 $\to$ 큰 쪽}
\end{tabular}}\\[1.0em]
{\large GDA 에서 $\Sigma$ 의 비대각을 0 으로 두면 $=$ \kb{가우시안 나이브 베이즈}}
\vfill
\note{나이브 베이즈와 GDA는 같은 가족, 생성 모델입니다. 둘 다 종마다 특징이 어떻게 나오는지 배우고, 사전확률을 곱해 베이즈 정리로 뒤집어 큰 쪽을 고릅니다. 다른 점은 특징의 종류와 모양입니다. 나이브 베이즈는 단어가 있다 없다 같은 0과 1 특징을 단어마다 베르누이로 배우고, 단어끼리 독립이라고 가정합니다. GDA는 실수 특징을 가우시안으로 배우고, 공분산으로 특징끼리 함께 움직이는 것까지 담습니다. 거꾸로 GDA에서 공분산의 비대각을 0으로, 즉 특징끼리 독립이라고 두면 가우시안 나이브 베이즈가 됩니다.}
\end{frame}

\section{연습}
% --- 문제풀이
\begin{frame}{문제풀이}
\vfill
\begin{enumerate}
\item 두 종의 평균이 4와 6, 공유 표준편차 1, 사전확률 반반. 경계는? \; $x = 5$ 에서 확률은?\vskip0.6em
\item 1번에서 $w_1 = (\mu_1 - \mu_0)/\sigma^2$ 은? 표준편차가 2 로 커지면 $w_1$ 은? 곡선은 어떻게 바뀌나?\vskip0.6em
\item 1번에서 종 1 이 90\% 인 밭이라면 경계는 왼쪽? 오른쪽? \; ($w_0$ 에 $\log\frac{\phi}{1-\phi}$ 가 더해진다)\vskip0.6em
\item 학습 데이터가 30개뿐이고 특징이 가우시안에 가깝다. GDA 와 로지스틱회귀 중 먼저 해 볼 것은?\vskip0.6em
\item $\Sigma = \begin{bmatrix}1 & -0.8\\ -0.8 & 1\end{bmatrix}$ 인 종의 타원은 어느 쪽으로 기울었나?
\end{enumerate}
\vfill
\note{1번 답. 표준편차가 같고 사전확률이 반반이면 경계는 두 평균의 한가운데, 5입니다. 경계에서는 두 높이가 같으니 확률은 0.5입니다. 2번 답. w1은 6 빼기 4를 1로 나눈 2입니다. 표준편차가 2가 되면 분산이 4라서 w1은 0.5, 기울기가 작아져서 S 자가 완만해집니다. 두 종이 많이 겹치면 확신이 줄어드는 겁니다. 3번 답. 로그 9, 약 2.2가 w0에 더해지니 같은 x에서 종 1일 확률이 커지고, 경계는 왼쪽, 종 0 쪽으로 옮겨 갑니다. 흔한 종 쪽으로 더 많이 판정하는 것입니다. 4번 답. 데이터가 적고 가정이 맞을 것 같으니 GDA를 먼저 해 볼 만합니다. 5번 답. 비대각이 음수라 한쪽이 크면 다른 쪽이 작은, 왼쪽 위에서 오른쪽 아래로 내려가는 타원입니다.}
\end{frame}

\section{정리}
% --- 정리
\begin{frame}{정리}
\vfill
\centering\large\begin{tabular}{l}
생성 모델 \quad 종마다 $P(x \mid y)$ 와 $P(y)$ 를 배우고 베이즈로 뒤집는다\\[0.5em]
GDA \quad $P(x \mid y)$ $=$ 가우시안, 학습 $=$ 비율 · 평균 · 공분산 재기\\[0.5em]
같은 $\Sigma$ \quad $x^2$ 가 지워져 $P(y \mid x) = \mathrm{sigmoid}(w_0 + \mathbf{w}\cdot\mathbf{x})$, 경계는 직선\\[0.5em]
다른 $\Sigma_k$ \quad 곡선 경계 (QDA)\\[0.5em]
생성 vs 판별 \quad 가정이 맞으면 적은 데이터에 강하다, 틀리면 로지스틱이 튼튼\\[0.5em]
붓꽃 \quad 4.9 cm: GDA 49\% · 로지스틱 58\%, 2차원 정확도 94\%
\end{tabular}
\vfill
\note{정리합니다. 생성 모델은 종마다 특징이 어떻게 나오는지와 종의 비율을 배우고 베이즈 정리로 뒤집어 분류합니다. GDA는 특징이 가우시안이라고 두고, 학습은 비율, 평균, 공분산을 재는 것이 전부입니다. 공분산을 같이 쓰면 x제곱이 지워져서 사후확률이 로지스틱 꼴이 되고 경계는 직선입니다. 종마다 다르게 쓰면 곡선 경계입니다. 가정이 맞으면 적은 데이터에 강하고, 틀리면 로지스틱회귀가 튼튼합니다. 다음 시간에는 확률 모델을 버리고 경계의 여유, 마진으로 가는 SVM을 봅니다.}
\end{frame}

% --- 참고 문헌
\begin{frame}{참고 문헌}
\vfill
\small
\begin{itemize}
\item Ng, A., \& Ma, T. CS229 Lecture Notes, Part IV: Generative Learning Algorithms. Stanford. \paperlink{https://cs229.stanford.edu/notes2021fall/cs229-notes2.pdf}
\item Ng, A. (2018). Lecture 5: GDA \& Naive Bayes, CS229 Autumn 2018. \paperlink{https://www.youtube.com/playlist?list=PLoROMvodv4rMiGQp3WXShtMGgzqpfVfbU}
\item Ng, A. Y., \& Jordan, M. I. (2001). On discriminative vs. generative classifiers: A comparison of logistic regression and naive Bayes. \textit{NeurIPS}.
\item Fisher, R. A. (1936). The use of multiple measurements in taxonomic problems. \textit{Annals of Eugenics}, 7(2).
\end{itemize}
\vfill
\note{오늘 내용은 스탠퍼드 CS229 강의 노트의 4부, 생성 학습 알고리즘을 바탕으로 했고, 앤드류 응 2018년 강의 5강에서 같은 내용을 볼 수 있습니다. 생성 모델과 판별 모델의 비교는 앤드류 응과 마이클 조던의 2001년 논문이 유명합니다. 붓꽃 데이터는 피셔의 1936년 논문입니다.}
\end{frame}
\end{document}
